\begin{abstractpage}

\begin{abstract}{romanian}
Lucrarea de față prezintă proiectarea și dezvoltarea aplicației web CourseHub, o
platformă care conectează profesori și elevi într-un spațiu comun de învățare. Obiectivul
principal este construirea unui sistem în care un utilizator își poate crea un profil ca
profesor sau ca elev, poate căuta un partener potrivit după materie, nivel și rol, poate
iniția o colaborare printr-o cerere de colaborare și poate comunica ulterior în timp real,
inclusiv prin partajarea de fișiere și materiale didactice. Platforma a fost dezvoltată
folosind ASP.NET Core 8, Entity Framework Core și SQL Server, cu comunicare
în timp real realizată prin SignalR. Printre funcționalitățile implementate se numără
autentificarea securizată prin cod OTP trimis pe email, un sistem de roluri bazat pe
Identity Roles cu panou de administrare dedicat, protecția formularului de înregistrare
prin reCAPTCHA, încărcarea de avatare și fișiere validate atât din punct de vedere al
tipului, cât și al dimensiunii, precum și o interfață disponibilă integral în limbile română
și engleză. Rezultatele obținute arată o platformă funcțională, cu accent pe securitate,
experiența utilizatorului și posibilitatea de moderare a conținutului, oferind în același
timp un punct de plecare solid pentru extinderi ulterioare.
\end{abstract}

\begin{abstract}{english}
This thesis presents the design and development of the CourseHub web application, a
platform that connects teachers and learners within a shared learning space. The main
objective is to build a system in which a user can create a profile as either a teacher
or a learner, search for a suitable partner by subject, level and role, initiate a collaboration
through a collaboration request, and subsequently communicate in real time, including through
file and teaching material sharing. The platform was developed using ASP.NET Core 8,
Entity Framework Core and SQL Server, with real-time communication implemented
through SignalR. Among the implemented features are secure authentication through
an email OTP code, a role-based system built on Identity Roles with a dedicated admin
panel, registration form protection through reCAPTCHA, avatar and file uploads
validated by both type and size, and a user interface fully available in both Romanian
and English. The results show a functional platform with an emphasis on security,
user experience and content moderation capabilities, while also providing a solid
foundation for future extensions.
\end{abstract}

\end{abstractpage}
